% pref.tex  See http://joshua.smcvt.edu/linearalgebra
{\setlength{\parskip}{.7ex}  % note the group-starting open curly
% \bigskip
% \vspace*{1.25in plus .2in minus .1in}
\chapter*{Prakata}
Buku ini membantu mahasiswa menguasai materi kuliah pertama Aljabar Linear
yang lazim pada jenjang sarjana di Amerika Serikat.

Materinya bersifat baku karena topik yang dibahas meliputi
reduksi Gauss,
ruang vektor, pemetaan linear,
determinan, serta nilai eigen dan vektor eigen.
Sasaran pembacanya pun lazim:
mahasiswa tahun kedua atau ketiga, biasanya dengan bekal
sekurang-kurangnya satu semester kalkulus.
Buku ini membantu mahasiswa melalui pendekatan yang bertahap\Dash
penyajiannya menekankan motivasi dan kewajaran gagasan,
dengan menggunakan banyak contoh.

Pendekatan bertahap inilah keunggulan utama buku ini,
jadi akan saya uraikan lebih lanjut.
Mata kuliah pada awal program matematika
lebih sedikit berfokus pada teori dan lebih banyak pada perhitungan.
Mata kuliah lanjutan
menuntut kematangan matematis:~kemampuan mengikuti berbagai
jenis argumen,
keakraban dengan
tema-tema yang mendasari banyak penyelidikan matematis, seperti
fakta dasar tentang himpunan dan fungsi,
serta kemampuan membaca dan berpikir secara mandiri sampai taraf tertentu.
Sebagian program memiliki mata kuliah tersendiri untuk mengembangkan kematangan ini,
tetapi bagaimanapun juga mata kuliah Aljabar Linear
merupakan tempat yang ideal untuk menjalani peralihan tersebut.
Mata kuliah ini hadir cukup awal sehingga kemajuan yang dicapai di sini
bermanfaat pada tahap-tahap selanjutnya,
namun juga cukup lanjut sehingga kelasnya berisi
mahasiswa yang sungguh-sungguh menekuni matematika.
Materinya mudah dijangkau, padu, dan anggun.
Contohnya pun melimpah.

Membantu pembaca menjalani peralihan ini
menuntut agar matematikanya diperlakukan dengan sungguh-sungguh.
Semua hasil dalam buku ini dibuktikan.
Di sisi lain, kita tidak dapat
menganggap mahasiswa sudah mencapai kematangan tersebut.
Karena itu,
berbeda dari buku-buku yang lebih lanjut,
buku ini dipenuhi ilustrasi teori,
yang sering kali diuraikan dengan cukup terperinci.

Sebagian buku yang ditujukan kepada pembaca yang belum matang secara matematis
dimulai dengan perkalian matriks dan determinan.
Kemudian, ketika
ruang vektor dan pemetaan linear akhirnya muncul
serta definisi dan pembuktian mulai digunakan, perubahan yang mendadak itu
membuat mahasiswa seketika tersendat.
% They've been sent the wrong signals.
Walaupun buku ini dimulai dengan
reduksi linear, sejak awal
kita melakukan lebih dari sekadar menghitung.
Bab pertama
memuat pembuktian, misalnya pembuktian bahwa reduksi linear menghasilkan
himpunan solusi yang benar dan lengkap.
Dengan itu sebagai motivasi,
bab kedua membahas ruang vektor atas bilangan real.
Dalam jadwal di bawah, pembahasan tersebut dimulai pada awal minggu ketiga.

Kemajuan terbesar seorang mahasiswa dalam matematika diperoleh dengan mengerjakan latihan.
Kumpulan soal dimulai dengan
pemeriksaan rutin dan berlanjut hingga pembuktian yang cukup mendalam.
% Since instructors often assign about a dozen exercises
Saya berusaha menyediakan sekitar dua lusin soal dalam setiap kumpulan,
sehingga tersedia cukup banyak pilihan.
Secara khusus, terdapat cukup banyak soal dengan tingkat kesulitan sedang
yang menantang pembelajar, tetapi tidak secara berlebihan.
Pada tingkat tertinggi, beberapa soal berupa teka-teki
yang diambil dari berbagai jurnal, kompetisi, atau
kumpulan soal; soal-soal ini
ditandai dengan
`\puzzlemark'
(demi keseruannya, saya berusaha mempertahankan redaksi aslinya).

Jadi, sebagaimana bagian lain buku ini,
latihan-latihannya bertujuan membangun kemampuan sekaligus
membantu mahasiswa merasakan nikmatnya
\emph{mengerjakan} matematika.
Mahasiswa sepatutnya melihat bagaimana gagasan-gagasan itu muncul dan mampu
membayangkan diri mereka mengerjakan hal serupa.


%\vspace*{.5in}
\medskip
\noindent{\bf Penerapan.}
%\smallskip
% The point of view taken here, that students should think of 
% Linear Algebra as about vector spaces
% and linear maps, is not taken to the complete exclusion of others.
Penerapan dan komputasi merupakan aspek yang menarik sekaligus penting
dalam bidang ini.
Karena itu, setiap bab ditutup dengan sejumlah
topik dari kedua aspek tersebut.
Topik-topik itu memberi pembaca
gambaran awal, membahas peran Aljabar Linear,
menunjukkan bacaan lanjutan, dan menyajikan beberapa latihan.
Uraiannya cukup singkat sehingga pengajar dapat membahas satu topik
dalam satu pertemuan
atau menugaskannya sebagai proyek perseorangan maupun kelompok kecil.
Terlepas dari apakah topik-topik itu menjadi bagian formal mata kuliah,
topik tersebut membantu pembaca melihat sendiri bahwa Aljabar Linear adalah
perangkat yang wajib dimiliki seorang profesional.




\medskip
\noindent{\bf Ketersediaan.}
Buku ini Bebas.
Lihat \url{hefferon.net/linearalgebra}
untuk rincian lisensinya.
Laman tersebut juga menyediakan versi terbaru,
jawaban latihan, salindia Beamer, panduan laboratorium, materi tambahan,
dan sumber \LaTeX.
Buku ini juga tersedia dalam bentuk cetak
dari penerbit-penerbit umum dengan biaya yang sangat terjangkau.
Lihat laman web tersebut.



\medskip
\noindent{\bf Ucapan Terima Kasih.}
Salah satu pelajaran dari pengembangan
perangkat lunak ialah bahwa proyek yang rumit mengandung kekeliruan
dan memerlukan proses untuk memperbaikinya.
Saya berterima kasih atas laporan dari para pengajar maupun mahasiswa.
Secara berkala saya menerbitkan revisi dan mencantumkan penghargaan
dalam repositori buku untuk setiap laporan yang saya gunakan.
Informasi kontak terbaru saya tersedia di laman web.

Saya berterima kasih kepada Saint Michael's College
karena telah mendukung proyek ini selama bertahun-tahun, bahkan sebelum gagasan
sumber daya pendidikan terbuka dikenal luas.
% I also thank Gabriel S Santiago for the cover colors.
% I also thank Adobe Color~CC user \texttt{claflin61} for the cover colors.

Dan saya tidak akan pernah cukup berterima kasih kepada istri saya, Lynne,
atas dukungannya yang tak pernah surut.



\newcommand{\classday}[1]{\textsc{#1}}
\newcommand{\colwidth}{1.25in}

%\vspace*{.5in}
\medskip
% \noindent{\bf If you are reading this on your own.}
\noindent{\bf Saran.}
%\smallskip
%
Penekanan buku ini pada motivasi dan perkembangan bertahap,
serta ketersediaannya, membuatnya banyak digunakan untuk belajar mandiri.
Jika Anda belajar secara mandiri, itu patut dihargai; saya mengagumi ketekunan Anda.
Namun, beberapa saran berikut mungkin berguna.

Seorang pengajar berpengalaman mengetahui topik dan
tempo yang sesuai untuk kelasnya. Meskipun demikian, jadwal semester berikut
(yang dengan baik hati dibagikan oleh G~Ashline)
dapat membantu Anda merencanakan laju belajar yang masuk akal.
Jadwal ini mengandaikan bahwa Anda telah mempelajari materi
Bagian~Satu.II, yaitu dasar-dasar vektor.
% \begin{center}   % George Ashline's
%    \begin{tabular}{r|*{2}{p{\colwidth}}l}
%       \multicolumn{1}{r}{\textit{week}}  
%        &\textit{Monday}          
%        &\textit{Wednesday}            
%        &\textit{Friday}        \\ \hline
%        1    &One.I.1         &One.I.1, 2        &One.I.2, 3         \\
%        2    &One.I.3         &One.III.1          &One.III.2         \\
%        3    &Two.I.1         &Two.I.1, 2         &Two.I.2         \\
%        4    &Two.II.1         &Two.III.1         &Two.III.2         \\
%        5    &Two.III.2        &Two.III.2, 3         &Two.III.3        \\
%        6    &\classday{exam}   &Three.I.1         &Three.I.1       \\
%        7    &Three.I.2         &Three.I.2          &Three.II.1         \\
%        8    &Three.II.1        &Three.II.2          &Three.II.2          \\
%        9    &Three.III.1       &Three.III.2         &Three.IV.1, 2       \\
%       10    &Three.IV.2, 3   &Three.IV.4          &Three.V.1          \\
%       11    &Three.V.1       &Three.V.2            &Four.I.1         \\
%       12    &\classday{exam}  &Four.I.2            &Four.III.1       \\
%       13    &Five.II.1    &\multicolumn{2}{c}{\classday{--Thanksgiving break--}} \\
%       14    &Five.II.1, 2     &Five.II.2          &Five.II.3        
%    \end{tabular}
% \end{center}
\begin{center} 
   \begin{tabular}{r|l}
      \multicolumn{1}{r}{\textit{Minggu}}
       &\textit{Bagian yang dibahas}        \\ \hline
       1    &Satu.I.1--2                 \\
       2    &Satu.I.3, Satu.III.1--2      \\
       3    &Dua.I.1--2       \\
       4    &Dua.II.1, Dua.III.1--2         \\
       5    &Dua.III.2--3      \\
       6    &\classday{--Ujian--}, Tiga.I.1      \\
       7    &Tiga.I.2, Tiga.II.1         \\
       8    &Tiga.II.1--2          \\
       9    &Tiga.III.1--2,  Tiga.IV.1--2       \\
      10    &Tiga.IV.2--4, Tiga.V.1          \\
      11    &Tiga.V.1--2, Empat.I.1         \\
      12    &\classday{--Ujian--},  Empat.I.2, Empat.III.1       \\
      13    &Lima.II.1, \classday{--Libur Thanksgiving--} \\
      14    &Lima.II.1--3        
   \end{tabular}
\end{center}
Sebagai pengayaan, Anda dapat memilih satu atau dua topik tambahan yang menarik
dari panduan laboratorium atau dari bagian Topik di akhir setiap bab.
Saya menyukai Topik tentang
Paradoks Pemungutan Suara,
Geometri Pemetaan Linear, dan Osilator Terkopel.
% (When I teach with this schedule as a target, I often have room
% for a couple of days of extras.)
Anda akan memperoleh manfaat lebih besar dari topik-topik ini
jika memiliki akses ke perangkat lunak perhitungan seperti
\textit{Sage}, yang tersedia bebas
di \url{http://sagemath.org}.

Dalam daftar isi,
saya menandai beberapa subbagian sebagai opsional karena
sebagian pengajar mungkin melewatinya agar dapat meluangkan lebih banyak waktu untuk materi lain.

Perhatikan bahwa,
selain ujian di kelas,
mahasiswa pada mata kuliah di atas mengerjakan
kumpulan soal untuk dibawa pulang yang memuat pembuktian, misalnya memverifikasi
bahwa suatu himpunan merupakan ruang vektor.
Perhitungan itu penting, tetapi argumen pun demikian.

Saran utama saya: kerjakanlah banyak latihan.
Saya telah menandai sampel yang baik dengan \recommendationmark\ di margin.
Jangan hanya membaca jawabannya\Dash Anda harus
mencoba mengerjakan soal dan mungkin perlu bergulat dengannya.
Untuk setiap latihan, Anda harus membenarkan jawaban dengan perhitungan
atau pembuktian.
Ingatlah bahwa hanya sedikit orang yang dapat menulis pembuktian yang benar tanpa latihan;
carilah orang berpengetahuan yang dapat bekerja bersama Anda.

Terakhir, sebuah peringatan bagi semua mahasiswa, baik yang belajar mandiri maupun tidak:~saya
tidak dapat cukup menekankan bahwa
pernyataan, ``Saya memahami materinya, hanya saja
saya kesulitan mengerjakan soalnya''\spacefactor=1000\ %
menunjukkan suatu kesalahpahaman.
Seluruh tujuan gagasan-gagasan tersebut adalah agar kita mampu menggunakannya.
Kutipan di bawah mengungkapkan pandangan ini dengan sangat baik
(saya mengambil kebebasan untuk menatanya sebagai puisi).
Keduanya menangkap hakikat keindahan sekaligus kekuatan
matematika dan sains pada umumnya,
serta Aljabar Linear pada khususnya.

\medskip
\par\noindent\begin{tabular}[t]{@{}l@{}}
  \textit{Saya tidak mengenal cara yang lebih baik}                     \\
  \textit{\ daripada menjelaskan prinsip-prinsip yang menggugah} \\
  \textit{melalui contoh-contoh khusus yang dipilih dengan cermat.}                    \\
  \hspace*{1in}\textit{--Stephen Jay Gould}
\end{tabular}

\medskip
\par\noindent
\begin{tabular}[t]{@{}l@{}}   
\textit{Jika Anda sungguh ingin belajar}                     \\
   \textit{\ Anda harus menaiki mesin itu}  \\ 
   \textit{\ dan mengenali segala muslihatnya} \\
   \textit{melalui percobaan nyata.}                    \\
   \hspace*{1in}\textit{--Wilbur Wright}
\end{tabular}

% \medskip
% \par\noindent
% \begin{tabular}[t]{@{}l@{}}   
% \textit{In the particular}                     \\
%    \textit{\ is contained the universal.}  \\ 
%    \hspace*{1in}\textit{--James Joyce}
% \end{tabular}

\vspace*{3ex}
\par\ \hfill\begin{tabular}[t]{@{}l@{}}
                       Jim Hef{}feron            \\
                       Matematika dan Statistika \\ 
                       Saint Michael's College \\ 
                       Colchester, Vermont\ USA 05439  \\     
                       \url{hefferon.net/linearalgebra} \\
                       % Date used in both book and answers for current version
12 Okt 2021
  \\
                       Edisi keempat, cetakan kedua
                    \end{tabular}

\vspace{3ex plus 1fill}
\par\noindent\textit{Catatan Penulis.}
Menciptakan latihan yang baik, yang mencerahkan sekaligus menguji,
merupakan tindakan kreatif dan kerja keras.
Penciptanya layak mendapat pengakuan.
Namun, secara tradisional buku-buku tidak mencantumkan atribusi
untuk soal.
Dalam buku ini saya mengubah kebiasaan tersebut jika saya yakin akan sumbernya.
Saya dengan senang hati menerima bantuan siapa pun agar saya dapat memberikan
atribusi yang tepat kepada soal-soal lainnya.
} % ends the open curly for the parskip from the top of this file
